\documentclass{situationist-preprint}

\title{The Conditions of Justified Civil Disobedience}
\msid{ETH-17-a}
\msfamily{ETH-17}
\msdiscipline{Normative and Applied Ethics}
\mstype{Principal manuscript}
\mskeywords{civil disobedience; duty to obey the law; fair play; necessity; publicity; punishment; rescue}

\begin{document}
\maketitle

\begin{abstract}
May one break a just law to prevent harm to others? If there is a duty to obey the law, it must be weighed against the reasons to prevent the harm, and the classical accounts of civil disobedience add that the breach must be public and its penalty accepted. These further conditions are usually defended as marks of the disobedient's sincerity or as means of persuasion, which leaves it unclear why they should bear on permissibility. We argue that one may break the law to prevent harm when the harm is grave, the breach minimal and public, and the legal consequences accepted. The argument starts from the view that the duty to obey just laws is weak, defeasible and grounded in fairness. We show that each condition answers to a feature of that ground. Gravity of harm and minimality of breach ensure that the reasons for the breach outweigh the duty; publicity and the acceptance of consequences ensure that the breach is not free-riding, since the agent neither conceals her exemption from those who bear the scheme's burdens nor withdraws from the burden the scheme imposes on those who break its rules. The conditions are jointly sufficient.
\end{abstract}

\section{The problem}

In the \emph{Crito}, Socrates refuses to escape from prison, arguing that by remaining in Athens he had agreed to abide by its laws, and that to break them now would be to injure the city that had raised him \citep[50a--54d]{PlatoCrito}. Thoreau, refusing to pay a tax that supported slavery and war, argued the opposite: that a just person's first allegiance is to what is right, and that a government which makes one an agent of injustice should be disobeyed \citep{Thoreau1849}. King, writing from jail in Birmingham, held both thoughts together: one who breaks an unjust law must do so openly, lovingly, and with a willingness to accept the penalty, and in doing so he expresses the highest respect for law \citep{King1964}.

Analytic philosophy took up the question in two strands. The first concerns the duty to obey the law. \citet{Hart1955} and \citet{Rawls1964} grounded it in fairness: those who benefit from a cooperative scheme of rules owe it to the others who bear its burdens to bear their share. \citet{Smith1973} and \citet{Simmons1979} argued that there is no general duty to obey, even a prima facie one. The second strand concerns civil disobedience. \citet[§§55--59]{Rawls1971} defined it as a public, nonviolent, conscientious, political act contrary to law, done with the aim of bringing about a change in law or policy, and held that the disobedient's acceptance of the legal consequences expresses fidelity to law. \citet[ch.~8]{Dworkin1977} argued that acceptance of punishment is not required, and that a society should often not prosecute. \citet[ch.~14]{Raz1979} argued that in a liberal state there is no right to civil disobedience, though disobedience may be justified. More recently, \citet{Lefkowitz2007} argued for a moral right to civil disobedience, \citet{Brownlee2012} defended conscientious disobedience as a communicative act whose constraints derive from that function, and \citet{Delmas2018} argued that there is sometimes a duty to resist, and that the resistance need not be civil.

Our question is narrower than the question of civil disobedience in Rawls's sense. It concerns breaches of law whose purpose is to prevent harm to others: breaking into a locked car to free a child in danger, trespassing to stop an unlawful discharge of poison into a river, disclosing protected information to prevent an impending catastrophe. Such breaches need not aim at a change in law. The law's own defences of necessity sometimes cover them, but often do not, because the harm is not imminent enough, or because the legislature has foreseen the situation and decided against the agent. Our question is when such a breach is permissible.

We answer: when the harm is grave, the breach minimal and public, and the legal consequences accepted. The answer is not new in its elements. What we add is a derivation of the conditions from the ground of the duty that the breach infringes. If the duty to obey just laws is a weak, defeasible duty of fairness, as POL-03 argues, then each condition answers to a feature of that duty, and together they are sufficient for the breach to be permissible.

\section{Framework}

\begin{definition}[Just law]\label{d:just}
A \emph{just law} is a rule of a legal system that is reasonably just as a whole, which is not itself seriously unjust, and whose general observance is part of a cooperative scheme from which those subject to it benefit.
\end{definition}

\begin{assumption}[The duty to obey]\label{a:duty}
There is a weak, defeasible duty to obey just laws, grounded in fairness: those who benefit from a cooperative scheme of rules owe it to those who bear the scheme's burdens to bear their share.
\end{assumption}

Assumption~\ref{a:duty} is the conclusion of POL-03, and it is the version of the fair-play account defended by \citet{Hart1955} and \citet{Rawls1964}. It is weak: it can be outweighed by reasons of moderate strength. It is defeasible: it lapses where its conditions fail. And it is grounded in fairness, which means that what is wrong with a breach, when it is wrong, is that the agent takes for herself a benefit the scheme provides, the general compliance of others, without bearing the burden that makes it possible.

\begin{definition}[The four conditions]\label{d:conditions}
A breach of law by an agent to prevent harm to others satisfies the \emph{four conditions} when:
\begin{enumerate}[label=(\roman*)]
\item \emph{gravity}: the harm is a serious setback to the basic interests of persons, such as their life, bodily integrity or liberty, and it cannot be prevented in time by lawful means available to the agent;
\item \emph{minimality}: the breach is no greater, in the burden it imposes on others and in the rules it infringes, than is needed to prevent the harm;
\item \emph{publicity}: the agent acts openly, or makes the breach and her reasons for it known to the relevant authorities promptly after acting;
\item \emph{acceptance}: the agent submits to the legal process and to the penalties it imposes for the breach, and does not evade them.
\end{enumerate}
\end{definition}

\begin{principle}[Free-riding]\label{p:freeride}
A breach of a just law is unfair, in the sense relevant to Assumption~\ref{a:duty}, if and only if the agent either conceals her exemption from those who bear the burdens of the scheme, or withdraws from the burden that the scheme attaches to its breach.
\end{principle}

\begin{principle}[Weighing]\label{p:weigh}
A weak duty is outweighed by the reasons to prevent a grave harm to others, provided that the infringement of the duty is no greater than is needed to prevent the harm.
\end{principle}

Principle~\ref{p:freeride} identifies what the fair-play account condemns. The free-rider takes the benefit of others' compliance while declining to comply, and two things make this unfair. The first is that he relies on others not knowing, or not doing likewise: his exemption is a private advantage, which he could not openly claim on terms that others would accept. The second is that he declines the burden that the scheme attaches to non-compliance, which is the scheme's way of making non-compliance costly and so keeping the burdens equal. An agent who breaks the law openly, on grounds she states and that would equally license anyone in her situation, and who accepts the penalty, does neither of these things. Principle~\ref{p:weigh} is the ordinary principle of proportion: the weaker a duty, the less it takes to outweigh it, but an infringement in excess of what the outweighing reason requires is not justified by it.

\section{The argument}

\begin{lemma}[No free-riding]\label{l:fair}
A breach that satisfies publicity and acceptance is not unfair in the sense of Principle~\ref{p:freeride}.
\end{lemma}

\begin{proof}
By publicity, the agent does not conceal her breach or her reasons; she claims her exemption openly, on grounds that those who bear the scheme's burdens can examine and that would apply to anyone in her situation. By acceptance, she bears the burden that the scheme attaches to breach. Neither disjunct of Principle~\ref{p:freeride} is satisfied, and the breach is not unfair.
\end{proof}

\begin{lemma}[Outweighing]\label{l:outweigh}
A breach that satisfies gravity and minimality is supported by reasons that outweigh the duty to obey.
\end{lemma}

\begin{proof}
By Assumption~\ref{a:duty}, the duty to obey is weak. By gravity, the breach prevents a grave harm to others that lawful means cannot prevent in time. By minimality, the infringement of the duty is no greater than is needed to prevent it. By Principle~\ref{p:weigh}, the reasons for the breach outweigh the duty.
\end{proof}

\begin{theorem}\label{t:main}
A breach of a just law to prevent harm to others is permissible if it satisfies the four conditions.
\end{theorem}

\begin{proof}
By Lemma~\ref{l:fair}, the breach does not involve the unfairness that grounds the duty to obey; and by Lemma~\ref{l:outweigh}, whatever force the duty retains is outweighed by the reasons to prevent the harm. No other duty has been assumed to bear against the breach. The breach is therefore permissible.
\end{proof}

The two lemmas do different work, and the theorem needs both. Gravity and minimality alone would make the breach supported on balance, but they would leave it a case of free-riding, and an agent who breaks the law in secret to prevent harm and then evades the consequences takes for herself an exemption she does not offer to others. Publicity and acceptance alone would remove the unfairness, but they would not make the breach worth its cost: a public breach with accepted penalties, made to prevent a trivial harm, still weakens a just scheme by example and imposes costs on others for no adequate reason. When all four are satisfied, the breach is both justified on balance and free of the specific wrong that the duty to obey protects against.

\begin{remark}\label{r:remainder}
On the account of ETH-15, the agent who meets the four conditions does not face a dilemma between law and rescue, since one option is uniquely permissible. The duty to obey is nonetheless infringed, and the acceptance of the legal consequences is the form taken, in this case, by the residual duties that infringed requirements leave behind. ETH-27 gives the same structure for agents who dirty their hands, where the residue takes the form of acknowledgment and amends.
\end{remark}

\begin{remark}\label{r:required}
Theorem~\ref{t:main} concerns permission. Whether the breach is also required depends on its cost to the agent. ETH-18 argues that there is a duty to rescue up to moderate cost, beyond which rescue is supererogatory. A breach whose accepted legal consequences are modest may fall within that duty; one that exposes the agent to imprisonment will usually exceed it, and is then permissible and admirable but not required.
\end{remark}

\section{Cases}

\paragraph{The child in the car.} A passerby sees a young child locked in a car on a hot day, in visible distress, with no adult nearby. She breaks a window, takes the child out, calls the emergency services and waits. The harm is grave and imminent; breaking one window is the least damage that will open the car; she acts openly and waits for the authorities; and she is ready to answer for the window. She satisfies the four conditions. In many jurisdictions the law's own necessity defence would excuse her, which confirms the verdict rather than making it redundant: the defence is the law's recognition of a class of breaches that the conditions independently justify.

\paragraph{The river.} An engineer learns that her firm is discharging a toxic effluent that will poison a town's water supply within days; regulators have been told and have not acted. She enters a locked control room, closes the outflow valve, and informs the police and the press of what she has done and why. The harm is grave and lawful means have failed in time. Closing a valve is minimal, whereas sabotaging the plant would not be. She acts publicly and accepts prosecution for trespass. She satisfies the four conditions.

\paragraph{The disclosure.} An official discloses classified documents showing that a policy is causing grave and continuing harm, after internal channels have failed. If he discloses only what is needed to reveal the harm, rather than everything he has access to, and surrenders to the authorities, as Daniel Ellsberg did in 1971 after the publication of the Pentagon Papers, he satisfies the four conditions. An official who discloses indiscriminately fails minimality; one who conceals his identity, or flees to avoid trial, fails publicity or acceptance. The theorem does not say that these officials act impermissibly; it says that its sufficient conditions are not met, and their cases must be assessed on other grounds.

\paragraph{Socrates.} Socrates was not preventing harm by escaping, and he did not consider his sentence a just application of a just law. His refusal to escape belongs to a different question, whether one must submit to an unjust verdict reached by a just procedure, which the companion manuscript ETH-17-b takes up. What the present account explains is why acceptance has the weight Socrates gave it. Acceptance is not a sign of obedience; it is what keeps a justified breach from being a private exemption.

\section{Objections}

\begin{objection}[Acceptance is not required]
Dworkin argued that the disobedient need not accept punishment, and that a just society should often refrain from punishing \citep[ch.~8]{Dworkin1977}. Requiring acceptance makes the permissibility of a good act depend on a willingness to suffer for it.
\end{objection}

\begin{reply}
Theorem~\ref{t:main} gives sufficient conditions, and we agree that a society may often rightly refrain from prosecuting those who meet them. What acceptance requires is that the agent not evade the legal process, not that the process punish her. If it declines to, she has borne what the scheme imposes, which is nothing. If it does punish her, her submission is what makes her breach consistent with the fairness she owes to those who comply; an agent who breaks a just law, however good her reason, and then escapes its sanction, holds herself exempt from a burden others bear.
\end{reply}

\begin{objection}[Publicity defeats rescue]
Some rescues succeed only if they are secret. Those who hid fugitives from slavery, or Jews from the Nazis, could not have announced what they were doing.
\end{objection}

\begin{reply}
These were breaches of seriously unjust laws, and by Definition~\ref{d:just} and Assumption~\ref{a:duty} the duty to obey does not extend to such laws: it is a duty to obey just laws, and where its conditions fail, it lapses. The rescuers did not need to meet the four conditions, which are conditions for justifying the breach of a just law. The objection confirms the account, since it shows that the demand for publicity disappears exactly where the fairness ground of the duty disappears.
\end{reply}

\begin{objection}[The laws will be destroyed]
The laws in the \emph{Crito} ask Socrates whether a city can survive in which private persons set aside its judgements. A permission to break just laws whenever one judges the harm grave is a permission whose general exercise would undermine the rule of law.
\end{objection}

\begin{reply}
The general exercise of the permission in Theorem~\ref{t:main} is the general exercise of breaches that are grave, minimal, public and accepted. A legal system can absorb such breaches, since each is open to examination, each is answered by the system's own process, and their number is limited by the gravity condition. What threatens the rule of law is the secret breach and the evaded sanction, and these are exactly what publicity and acceptance exclude. The Crito's argument, applied to the permission we defend, supports its conditions rather than its denial.
\end{reply}

\begin{objection}[No duty to obey]
If \citet{Smith1973} and \citet{Simmons1979} are right that there is no general duty to obey the law, the four conditions are unnecessary, and only the gravity of the harm matters.
\end{objection}

\begin{reply}
The argument assumes POL-03, and so assumes a duty to obey just laws. If there were no such duty, Lemma~\ref{l:fair} would be idle, and Theorem~\ref{t:main} would remain true, since its conditions would be more than sufficient. The philosophical anarchist must still explain what is objectionable about the secret breach whose consequences are evaded, when the harm prevented is grave; the fair-play account explains it, and to that extent the conditions give independent support to the duty they presuppose.
\end{reply}

\section{What the result does not establish}

The argument establishes that the four conditions are sufficient. It does not establish that each is necessary for every justified breach. Breaches of unjust laws are not governed by them, as the second objection shows, and there may be breaches of just laws that are justified without publicity because publicity would itself cause grave harm. Nor does the argument settle what counts as a grave harm at the margins, or how a society should respond to those who meet the conditions; we have suggested only that leniency is often appropriate. Finally, the argument concerns harm to others. Breaking the law to prevent harm to oneself raises different questions, since the reasons in play are of a different kind.

\section{Conclusion}

One may break a just law to prevent harm when the harm is grave, the breach minimal and public, and the legal consequences accepted. The duty to obey just laws is a weak duty of fairness. Gravity and minimality make the reasons for the breach outweigh it; publicity and acceptance ensure that the breach does not commit the wrong it guards against, since the agent neither hides her exemption nor escapes the burden the scheme attaches to it. Acceptance of the consequences is thus not a ritual of submission but the residue of a justified infringement, owed to those whose compliance the agent continues to rely on.

\begin{thebibliography}{99}\small
\bibitem[Brownlee(2012)]{Brownlee2012} Brownlee, K. (2012). \emph{Conscience and Conviction: The Case for Civil Disobedience}. Oxford University Press.
\bibitem[Delmas(2018)]{Delmas2018} Delmas, C. (2018). \emph{A Duty to Resist: When Disobedience Should Be Uncivil}. Oxford University Press.
\bibitem[Dworkin(1977)]{Dworkin1977} Dworkin, R. (1977). \emph{Taking Rights Seriously}. Cambridge, MA: Harvard University Press.
\bibitem[Hart(1955)]{Hart1955} Hart, H.~L.~A. (1955). Are there any natural rights? \emph{The Philosophical Review}, 64(2), 175--191.
\bibitem[King(1964)]{King1964} King, M.~L., Jr. (1964). \emph{Why We Can't Wait}. New York: Harper \& Row.
\bibitem[Lefkowitz(2007)]{Lefkowitz2007} Lefkowitz, D. (2007). On a moral right to civil disobedience. \emph{Ethics}, 117(2), 202--233.
\bibitem[Plato(Crito)]{PlatoCrito} Plato. \emph{Crito}. Cited by Stephanus page.
\bibitem[Rawls(1964)]{Rawls1964} Rawls, J. (1964). Legal obligation and the duty of fair play. In S. Hook (ed.), \emph{Law and Philosophy}. New York University Press.
\bibitem[Rawls(1971)]{Rawls1971} Rawls, J. (1971). \emph{A Theory of Justice}. Cambridge, MA: Harvard University Press.
\bibitem[Raz(1979)]{Raz1979} Raz, J. (1979). \emph{The Authority of Law: Essays on Law and Morality}. Oxford: Clarendon Press.
\bibitem[Simmons(1979)]{Simmons1979} Simmons, A.~J. (1979). \emph{Moral Principles and Political Obligations}. Princeton University Press.
\bibitem[Smith(1973)]{Smith1973} Smith, M.~B.~E. (1973). Is there a prima facie obligation to obey the law? \emph{The Yale Law Journal}, 82(5), 950--976.
\bibitem[Thoreau(1849)]{Thoreau1849} Thoreau, H.~D. (1849). Resistance to civil government. In E.~P. Peabody (ed.), \emph{Aesthetic Papers}. Boston.
\end{thebibliography}

\end{document}
